% GENERATED FILE. Edit canonical lessons, then run npm run generate:books.
% canonical-source-hash: fnv1a64:469d536837c30215
% canonical-lessons: HI-C32-shubh-sandhya, HI-S145-letter-sha, HI-C32-evening-register

\chapter{Good Evening}
\label{ch:hi-good-evening}
\hlchaptermodality{32}

\begin{chapteropening}
\textbf{By the end of this chapter:} \emph{I can wish someone good evening in Hindi and choose between \hi{शुभ} \hi{संध्या} and \hi{नमस्ते} by register.}

Pick \hi{शुभ} \hi{संध्या} for a written evening card and \hi{नमस्ते} for a casual 7 p.m. hello, treating the split as a tendency, not a ban.
\end{chapteropening}

\section[shubh sandhyā / sāñjh]{\hi{शुभ} \hi{संध्या} (shubh sandhyā) — \textquotedblleft{}good evening,\textquotedblright{} a twilight word with two lives}

\label{lesson:HI-C32-shubh-sandhya}



\begin{quote}
\hi{शाम} was Persian. This greeting reaches for something else again — a native Sanskrit word for the exact moment day and night meet.
\end{quote}

\subsection*{You'll want to know — \hi{संध्या} — \textquotedblleft{}the junction of day and night\textquotedblright{}}
\textbf{\hi{संध्या}} (\textbf{sandhyā}) — \textquotedblleft{}\textbf{evening, twilight}\textquotedblright{} — is \textbf{native Sanskrit}, from \textbf{\hi{सम्}} (\emph{sam-}, \textquotedblleft{}together\textquotedblright{}) + \textbf{\hi{धा}} (\emph{dhā}, \textquotedblleft{}to hold, to place\textquotedblright{}), literally \textquotedblleft{}\textbf{holding together}\textquotedblright{} — the \textbf{junction} where day and night meet. This is a genuinely different word from \textbf{\hi{शाम}} (\emph{shaam}, HI-C29's Persian loan): Hindi's everyday evening noun and its \textquotedblleft{}good evening\textquotedblright{} greeting don't share a root, the same pattern already found for morning.

\subsection*{You'll want to know — \hi{साँझ} — a doublet, but not a simple \textquotedblleft{}casual\textquotedblright{} version}
\hi{संध्या} has its own Hindi \textbf{tadbhava} doublet, \textbf{\hi{साँझ}} (\textbf{sāñjh}), worn down through Prakrit the same way \textbf{\hi{रात्रि}} became \textbf{\hi{रात}} (HI-C26/27). But be careful not to force that exact parallel: unlike \emph{rātri}/\emph{raat}, where the tadbhava form clearly became the everyday word and the tatsama stayed formal, sources describe \textbf{\hi{साँझ}} as used across \textbf{both} literary/poetic \textbf{and} everyday communication — not simply \textquotedblleft{}the casual one.\textquotedblright{} The doublet exists; the clean formal/casual split doesn't carry over as neatly this time.

\subsection*{Your turn}
\begin{itemize}
\raggedright
  \item \emph{Say it:} \textquotedblleft{}shubh sandhyā\textquotedblright{} — \textquotedblleft{}good evening,\textquotedblright{} su + sandhyā
  \item \emph{Say it:} sandhyā's own meaning — \textquotedblleft{}holding together,\textquotedblright{} the junction of day and night
  \item \emph{Contrast:} tatsama \emph{sandhyā} · evolved doublet \emph{sāñjh}
\end{itemize}

\begin{tcolorbox}[breakable,colback=teal!4,colframe=teal!35!black,title={Before you move on}]
What does \hi{संध्या} literally mean, and what is it built from? (\textquotedblleft{}\textbf{Holding together}\textquotedblright{} — \hi{सम्} + \hi{धा} — i.e. the junction of day and night.) Does \hi{साँझ} work the same way \hi{रात्रि}/\hi{रात}'s tatsama/tadbhava split did — one strictly formal, one strictly casual? (\textbf{No} — \hi{साँझ} is used across both literary and everyday registers, not confined to one.) Next: choose between the written greeting and the casual all-day greeting.
\end{tcolorbox}
\section[śa]{\hi{श} — the letter you have been thanking people with}

\label{lesson:HI-S145-letter-sha}



\begin{quote}
The hook from before: where did it sit, and what did it replace?

Today is a whole letter, and one you have been looking at since your first page.
\end{quote}

\subsection*{Script you'll notice: \hi{श}}
\textbf{\hi{श}} — \emph{śa}, the \emph{sh} of English \emph{ship}.

What it is made of:

\begin{itemize}
\raggedright
  \item a joined upper loop, descending outer curve, lower loop and diagonal tail
  \item a top-to-bottom right stem
  \item the top shirorekhā
\end{itemize}

The left-hand side is more work than most letters ask for, and it is \textbf{all one run} — the pen does not come up until the tail is finished.

\subsection*{Script: the third sibilant}
\noindent
\begin{tabularx}{\linewidth}{@{}>{\raggedright\arraybackslash}X>{\raggedright\arraybackslash}X@{}}
\toprule
\textbf{letter} & \textbf{sound} \\
\midrule
\hi{स} & \emph{sa} \\
\hi{ष} & \emph{ṣa} \\
\hi{श} & \emph{śa} \\
\bottomrule
\end{tabularx}

You could already draw two of the three. \textbf{This is the last one}, and with it the hissing letters are complete.

Words that have been yours from the start:

\begin{itemize}
\raggedright
  \item \textbf{\hi{शुक्रिया}} — thank you
  \item \textbf{\hi{शाम}} — evening
  \item \textbf{\hi{शनिवार}} — Saturday
\end{itemize}

The first of those was among the first words in the book, and the letter that opens it is only being drawn now. The opening chapters taught whole words, and the letters inside them went unnamed.

\subsection*{Writing: \hi{श}}
\begin{itemize}
\raggedright
  \item \textbf{1.} start at the upper loop's lower inner tip, sweep left and clockwise around the upper loop, descend the outer curve, curl around the lower loop, and continue down-right through the diagonal tail without lifting
  \item \textbf{2.} lift and draw the right stem top-to-bottom
  \item \textbf{3.} lift and draw the top shirorekhā left-to-right
\end{itemize}

\textbf{Pen lifts: 2.}

\begin{quote}
Verified three-stroke teaching form; everyday handwriting may join or simplify the body differently.
\end{quote}

\begin{quote}
This is one attested teaching order and not a national standard — handwriting here is taught with school-to-school variation. Source: Opiaterein, ‘Deva-\hi{श}-order.gif’, strokes 1–3, Wikimedia Commons, 10 May 2009.
\end{quote}

\subsection*{Your turn}
\begin{itemize}
\raggedright
  \item \emph{Look:} at these two words and find \hi{श} in each one
\end{itemize}

\begin{quote}
\hi{शुक्रिया}  ·  \hi{शाम}
\end{quote}

\begin{itemize}
\raggedright
  \item \emph{Trace it:} \hi{श} three times, saying \emph{śa} as you finish each one
  \item \emph{Say it:} how many pen lifts the whole letter took
\end{itemize}

\begin{tcolorbox}[breakable,colback=teal!4,colframe=teal!35!black,title={Before you move on}]
What sound does \textbf{\hi{श}} carry? (\emph{śa}.) Which three letters make the hissing set? Name the word of thanks it opens.
\end{tcolorbox}
\section[shubh sandhyā / namaste]{\hi{शुभ} \hi{संध्या} or \hi{नमस्ते}? — choose by setting}

\label{lesson:HI-C32-evening-register}



\begin{quote}
\emph{Shubh sandhyā} is structurally transparent. The final question is where Hindi speakers actually reach for it.
\end{quote}

\subsection*{You'll want to know — Written warmth and casual speech}
Like \textbf{\hi{सुप्रभात}}, \textbf{\hi{शुभ} \hi{संध्या}} appears heavily in \textbf{written and social greeting-card culture}: greeting poetry, WhatsApp messages, and evening well-wishes. It is less likely to open an ordinary mid-conversation exchange.

\textbf{\hi{नमस्ते}} works at any time of day, evenings included, and is more common in casual conversation than either \emph{shubh prabhāt} or \emph{shubh sandhyā}. The native language has divided the jobs: a Sanskrit \textquotedblleft{}good \_\_\_\textquotedblright{} phrase for formal or written warmth, and another native greeting for everyday speech.

This is not English displacing a Hindi greeting. It is Hindi choosing between two of its own expressions by register and medium.

\subsection*{Your turn}
\begin{itemize}
\raggedright
  \item \emph{Choose:} a WhatsApp evening card $\to$ \textbf{\hi{शुभ} \hi{संध्या}}
  \item \emph{Choose:} greeting someone casually at 7 p.m. $\to$ \textbf{\hi{नमस्ते}}
  \item \emph{Say it:} \textquotedblleft{}formal/written is a tendency, not an absolute ban\textquotedblright{}
\end{itemize}

\begin{tcolorbox}[breakable,colback=teal!4,colframe=teal!35!black,title={Before you move on}]
Where does \textbf{\hi{शुभ} \hi{संध्या}} especially thrive? (\textbf{Written and social greeting-card culture}.) Does \textbf{\hi{नमस्ते}} genuinely work in the evening? (\textbf{Yes} — it works at any time and is common casually.) Is the distinction Hindi versus English? (\textbf{No} — it is a native register/medium choice.)
\end{tcolorbox}
